% Philippica XIV (philippic-14) --- English translation
% Translated by Claude (Anthropic) from the Latin (Perseus).
% Style guide: STYLE.md.

\ciceroSection{1} If, senators, just as I have learned from the dispatches read out to us that the army of those most criminal enemies has been cut to pieces and put to flight, so I had learned that Decimus Brutus has already come forth from Mutina --- the thing we all most earnestly long for, and which we judge has followed from the victory that has been won --- then, since it was for his peril that we took up the war-cloak, for that same man's safety I should move without any hesitation that we return to our former dress. But until the thing the city awaits so eagerly has been brought us, it is enough to take joy in a very great and most glorious battle: reserve the return to civilian dress for the victory once it is complete.

\ciceroSection{2} And the completion of this war is the safety of Decimus Brutus. But what sort of motion is this --- that for today our dress be changed, and then tomorrow we come forth in the war-cloak again? No: once we have returned to the dress we crave and pray for, let us see to it that we keep it for ever. For this, surely, is not only shameful but not even pleasing to the immortal gods: to leave the altars we have approached in the toga and go off to take up the war-cloak.

\ciceroSection{3} And I observe, senators, that certain men favour this motion whose mind and design is this: that, seeing what a most glorious day that will be for Decimus Brutus --- the day on which we return to the toga for his safety --- they crave to snatch that reward from him, so that it may not be handed down to memory and to posterity that for the peril of a single citizen the Roman people took up the war-cloak, and for that same man's safety returned to the toga. Take that motive away: you will find no reason for so perverse a proposal. But you, senators, preserve your authority, stand by your resolve, hold fast in your memory what you have often shown: that the whole issue of this war turns on the life of one most brave and great man.

\ciceroSection{4} To free Decimus Brutus, envoys were sent --- the leading men of the state --- to warn that enemy and parricide to withdraw from Mutina; for the sake of saving that same Decimus Brutus, a consul, Aulus Hirtius, set out for the war by lot, whose feeble health the courage of his spirit and the hope of victory made strong; Caesar, when with an army raised by his own means he had first freed the commonwealth from these plagues, set out --- that no further crime might arise thereafter --- to free that same Brutus, and overcame some private grief by his love of country.

\ciceroSection{5} What else did Gaius Pansa do --- in holding levies, in raising money, in passing the gravest decrees of the Senate against Antony, in exhorting us, in summoning the Roman people to the cause of liberty --- but work to set Decimus Brutus free? And the Roman people in throng so demanded the safety of Decimus Brutus with one voice that they set it before not only their own advantage but even the necessities of their daily bread. And this we ought indeed to hope, senators, is either close at hand or already accomplished; but it is fitting that the fruit of our hope be reserved for the fact and the outcome, lest we seem either to have snatched at the gift of the immortal gods by our haste, or in our folly to have despised the power of fortune.

\ciceroSection{6} But since your show of feeling sufficiently declares what you think on this matter, I will come to the dispatches sent by the consuls and the propraetor --- once I have first said a few things that bear on the dispatches themselves. The swords of our legions and armies have been dipped, senators --- or rather drenched --- in two battles of the consuls and in a third of Caesar's. If that blood was the enemy's, the soldiers' devotion was complete; if it was the blood of citizens, an unspeakable crime. How long, then, shall the man who has outdone all enemies in wickedness go without the name of enemy? --- unless you would have the very sword-points of our soldiers tremble, in doubt whether they are sunk in a citizen or in an enemy. You decree a thanksgiving: you do not call him an enemy.

\ciceroSection{7} Welcome indeed to the immortal gods will our thanksgivings be, welcome the victims, when a multitude of citizens has been slain! ``Against the wicked,'' someone says, ``and the reckless'' --- for so that most distinguished man calls them: terms of abuse out of the lawsuits of the city, not the brands burned in by a war of extermination. They forge wills, I suppose, or evict their neighbours, or swindle young men; for it is men tainted with these vices and their like that custom is wont to call ``wicked'' or ``reckless.''

\ciceroSection{8} One man, the foulest of all bandits, wages a war past atonement against four consuls; the same man wages war against the Senate and the Roman people; on all men --- though he is himself collapsing under his own disasters --- he threatens ruin, devastation, torture, the rack. Dolabella's savage and monstrous deed, which no barbarian land could acknowledge as its own, he attests was done by his own counsel; and what he would have done in this city, had not Jupiter himself here driven him back from this temple and these walls, he made plain in the calamity of the people of Parma. Those most excellent and honourable men, bound above all to the authority of this order and the dignity of the Roman people, that thing of shame and prodigy Lucius Antonius destroyed by the cruellest of precedents --- Lucius Antonius, the marked hatred of all mankind, or, if the gods too hate the men they ought, of the gods as well.

\ciceroSection{9} My spirit shrinks back, senators, and dreads to tell what Lucius Antonius wrought upon the children and wives of the Parmenses. The very acts of shame the Antonii gladly submitted to, to their own disgrace, they now rejoice to have inflicted by force on others. But the violence they offered is a thing of ruin; the lust with which the life of the Antonii is besmeared, a thing of infamy. Is there anyone, then, who would not dare to call them enemies, when he confesses that by their wickedness the cruelty of the Carthaginians has been surpassed? For in what captured city was Hannibal ever so savage as Antonius in Parma surprised? --- unless perhaps, toward this colony, and toward the rest he is minded the same way against, he is not to be thought an enemy.

\ciceroSection{10} But if beyond any doubt he is the enemy of the colonies and the country towns, what, then, do you judge him to this city, which he has coveted to fill up the destitution of his banditry --- the city which his practised surveyor, the cunning Saxa, had already parcelled out with his ten-foot rod? Recall, by the immortal gods, senators, what we feared these last two days, when the most wicked rumours were scattered abroad by enemies within our own walls. Who could look on his children, who on his wife, without weeping? who on his home, his roof, his household gods? Every man was thinking of either a most hideous death or a pitiable flight. And do we hesitate to call those men enemies, from whom these terrors came? If anyone brings forward a heavier name, I shall gladly assent; with this common one I am barely content --- a lighter one I will not use.

\ciceroSection{11} And so, since we ought to decree a most just thanksgiving out of the dispatches that were read aloud, and since Servilius has so moved, I will outright increase the number of days --- the more so since they are to be decreed not for one leader but for three. But first I will do this: I will name as imperatores the men by whose courage, counsel, and good fortune we have been freed from the gravest perils of slavery and destruction. For to whom in these last twenty years has a thanksgiving been decreed without his being hailed imperator --- and that for the slightest of exploits, or most often for none at all? Wherefore either the man who spoke before me should not have moved a thanksgiving at all, or the customary, common-handed honour must be granted to those to whom even new and singular honours are owed.

\ciceroSection{12} Or, if a man had killed a thousand or two thousand Spaniards or Gauls or Thracians, would the Senate, by this custom that has grown so common, hail him imperator --- yet, now that so many legions have been cut down, so great a multitude of enemies slain (enemies, I say, however these enemies within would have it otherwise), shall we grant our most illustrious leaders the honour of thanksgivings and take from them the imperator's name? For with what honour, what gladness, what congratulation ought those very liberators of this city to enter this temple, when only yesterday, for their deeds, the Roman people carried me from my house up to the Capitol in ovation and all but in triumph, and from there brought me home again?

\ciceroSection{13} For that, in my judgement at least, is the just and true triumph: when to men who have served the commonwealth well a testimony is rendered by the consent of the whole state. For whether, amid the common joy of the Roman people, they were congratulating one man, that is a great judgement; or whether they were thanking one man, a greater still; or whether both, nothing more magnificent can be conceived. ``You, then --- of your own self?'' someone may say. Unwillingly, I grant; but the smart of injury makes me boastful, against my habit. Is it not enough that men ignorant of virtue render no gratitude to those who deserve well of them? Shall accusation and envy be hunted out even against those who fix all their cares on the safety of the commonwealth?

\ciceroSection{14} For you know that in these last days a talk ran most persistently: that on the Parilia --- which day is today --- I would come down to the Forum with the fasces. This, I suppose, was levelled at some gladiator or bandit or Catiline --- not at the man who brought it about that nothing of the kind could happen in the commonwealth. Or shall I, who crushed, overthrew, and laid low Catiline as he plotted these very things, myself have turned out, all at once, a Catiline? By what auspices was I, an augur, to take up those fasces? how long was I to keep them? to whom hand them on? Could anyone have been so wicked as to invent it, so frenzied as to believe it? Whence, then, came that suspicion --- or rather, whence that talk?

\ciceroSection{15} When, as you know, for these three or four days a grim report was seeping from Mutina, the impious citizens, swollen with joy and insolence, gathered themselves into one place --- to that senate-house, ill-starred through their own madness rather than the commonwealth's. There, as they entered upon plans for our slaughter and parcelled out among themselves who should seize the Capitol, who the Rostra, who the city gates, they reckoned that the whole state would come running to me. And that this might happen to my discredit, and at the peril of my life besides, they put about that rumour of the fasces; the fasces they meant to bring to me themselves. Then, once it had been done as if by my own will, an onslaught of hired men was being made ready against me, as against a tyrant; and from that the slaughter of you all would have followed. This affair has been brought into the open, senators; but in its own time the wellspring of the whole crime will be laid bare.

\ciceroSection{16} And so Publius Apuleius, tribune of the plebs --- the witness, the confidant, the helper of all my counsels and perils ever since my consulship --- could not endure his own grief at my distress. He held a very great public meeting, with the Roman people of one and the same mind. And in that meeting, when, out of our closest union and friendship, he wished to clear me of the suspicion of the fasces, the whole assembly with one voice declared that I had never thought anything about the commonwealth that was not the best. Two or three hours after this meeting was held, the most longed-for messengers and dispatches arrived: so that one and the same day not only freed me from a most unjust ill-will but exalted me with the Roman people's most crowded congratulation.

\ciceroSection{17} I have brought this in, senators, not so much to speak in my own defence --- for it would go ill with me if I were not cleared in your eyes without a defence --- as to warn certain men of too meagre and narrow a spirit to do what I have always done myself: to count the virtue of outstanding citizens a thing worthy of imitation, not of envy. Wide is the field in public life, as Marcus Crassus used wisely to say, and the course to glory lies open to many. Would that those leading men were alive who, after my consulship --- when I myself was yielding place to them --- saw me as a leader, and not unwillingly! But at this present time, in so great a dearth of steadfast and brave men of consular rank, with what grief do you suppose I am afflicted, when I see some ill-affected, others caring nothing at all, others holding too unsteadily to the cause they have taken up and trimming their opinion not always by the good of the commonwealth, but now by hope and now by fear?

\ciceroSection{18} But if anyone troubles himself over a contest for primacy --- which ought not to exist at all --- he acts most foolishly if he contends against virtue with vices; for as a race is won by running, so among brave men virtue is outdone only by virtue. Will you, if I hold the best views about the commonwealth, hold the worst yourself in order to beat me? Or, if you see good men flocking to me, will you invite the wicked to yourself? I would not have it so --- first for the commonwealth's sake, then for your own dignity's. But if it were primacy at stake --- which I have never sought --- what could I wish for more? For by bad opinions I cannot be outdone; by good ones perhaps I can --- and gladly.

\ciceroSection{19} That the Roman people sees these things, marks them, and judges them, certain men take amiss. But could it have been otherwise --- that men should not judge each according to his desert? For just as the Roman people judges most truly of the Senate as a whole that at no season of the commonwealth has this order been firmer or braver, so of each one of us --- and chiefly of those of us who deliver opinions from this place --- all men inquire, and are eager to hear what each has thought: and so they reckon each as they judge he has deserved. They hold it in memory that I, on the thirteenth day before the Kalends of January, was the first to recall our liberty;

\ciceroSection{20} that from the Kalends of January to this hour I have kept watch over the commonwealth; that my house and my ears have lain open day and night to the counsels and warnings of all men; that by my letters, my messages, my exhortations all men everywhere have been roused to the defence of the fatherland; that by my motions, from the Kalends of January, never were envoys sent to Antony; that he was always to me the enemy, this always the war; so that I, who at every time had been the advocate of a true peace, was the foe of this name of a ruinous peace; and the same of Publius ---

\ciceroSection{21} Ventidius: when others would have him a praetor, I always an enemy. Had the consuls been willing to put these motions of mine to a division, the very authority of the Senate would long ago have struck the arms from the hands of all those bandits. But what was not allowed then, senators, is now not only allowed but necessary: that those who are in fact enemies be branded so in words, and judged enemies by our votes.

\ciceroSection{22} Before now, when I had named the enemy and the war, once and again they struck my motion from the number of motions to be put: a thing that in this case can no longer be done. For it is upon the dispatches of the consuls Gaius Pansa and Aulus Hirtius, and of Gaius Caesar the propraetor, concerning honour to be paid the immortal gods, that we are delivering our opinions. The man who has just decreed a thanksgiving has, all unwittingly, judged them enemies; for never in a civil war has a thanksgiving been decreed. Decreed, do I say?

\ciceroSection{23} It was not even asked for in the victor's dispatches. Sulla as consul waged a civil war; he led his legions into the city, drove out whom he would, killed whom he could: of a thanksgiving, no mention. A grievous war, the Octavian, followed: no thanksgiving for Cinna the victor. Sulla as imperator avenged Cinna's victory: no thanksgiving was decreed by the Senate. To you yourself, Publius Servilius, did your colleague send any dispatch about that most calamitous battle of Pharsalus? did he wish you to move a thanksgiving? Surely he did not. Yet afterwards he sent word of Alexandria, of Pharnaces; but for the battle of Pharsalus he did not so much as hold a triumph. For that battle had carried off citizens with whom --- not living only, but even victorious --- the state might have stood safe and flourishing.

\ciceroSection{24} And the same had befallen the earlier civil wars. For to me as consul a thanksgiving was decreed --- though no arms had been taken up --- not for the slaughter of enemies but for the preservation of citizens, by a new and unheard-of kind of honour. Wherefore either the thanksgiving must be refused, though the commonwealth has been most gloriously served and our commanders ask it --- which has befallen no man but Aulus Gabinius --- or else, in decreeing the thanksgiving, you must of necessity judge enemies the men about whom you decree it. What, then, that man does in deed, I do also in word, when I call them imperatores: by that very name I judge to be enemies both those already wholly defeated and those who survive.

\ciceroSection{25} For how else should I rather name Pansa, though he holds the most ample title of honour? how name Hirtius? He is indeed a consul; but the one name is the Roman people's gift, the other is of valour and victory. And Caesar --- begotten for the commonwealth by the gods' grace --- shall I hesitate to call him imperator? --- the man who first turned aside Antony's monstrous and loathsome cruelty not only from our throats but even from our limbs and our inward parts. Of a single day, how many and how great the deeds of valour, immortal gods, there were!

\ciceroSection{26} For Pansa was the first of all to give battle and to close with Antony; a commander worthy of the Martian legion, a legion worthy of its commander. Had Pansa been able to curb its furious charge, the whole business would have been finished in a single battle. But when the legion, greedy for liberty, had burst all unbridled into the enemy's line, and Pansa himself was fighting in the front rank, he took two dangerous wounds and, borne off from the field, kept his life for the commonwealth. And I judge him not an imperator only but a most illustrious imperator, who, having pledged to satisfy the commonwealth by either death or victory, achieved the one --- may the immortal gods avert the omen of the other!

\ciceroSection{27} What shall I say of Hirtius? who, when he heard the news, led two legions out of the camp with incredible ardour and valour --- that Fourth legion which long ago, deserting Antony, joined itself to the Martian, and the Seventh, which, made up of veterans, proved in this battle that to those soldiers who had kept faith with Caesar's bounty the name of the Senate and people of Rome was dear. With these twenty cohorts, and no cavalry, Hirtius himself bearing the eagle of the Fourth legion --- than which we have heard of no fairer sight of a commander --- closed with Antony's three legions and his cavalry, and those wicked enemies, who hung over this temple of Jupiter Best and Greatest and the other temples of the immortal gods, over the roofs of the city, the liberty of the Roman people, our life and our blood, he laid low, routed, and slew, so that with a very few, hidden by night and stricken with terror, the chief and captain of the bandits took to flight. O happiest Sun himself, who, before he hid his face, with the corpses of the parricides strewn about, saw Antony fleeing with his few!

\ciceroSection{28} But will anyone hesitate to call Caesar imperator? His age, surely, will deter no man from this verdict, seeing that by his valour he has overtopped his age. And to me the services of Gaius Caesar have always seemed the greater the less they were to be demanded of such years. When we gave him a command, at that same moment we were conferring also the hope of that name; and when he had won it, he confirmed the authority of our decree by his own exploits. This young man, then, of the greatest spirit --- as Hirtius most truly writes --- guarded the camp of many legions with a few cohorts, and fought a successful battle. So, by the valour, the counsel, and the good fortune of three imperatores, in one day and in several places the commonwealth was preserved.

\ciceroSection{29} I move, therefore, in the name of those three, a thanksgiving of fifty days; the grounds I will embrace in the motion itself, in the most honorific words I can command. And it belongs to our good faith and our devotion to show the bravest of soldiers how mindful we are of them, and how grateful. Wherefore I move that our promises --- the rewards we pledged to give the legions when the war is finished --- be renewed by today's decree of the Senate; for it is right that the soldiers' honour, of such soldiers above all, be joined with it.

\ciceroSection{30} And would, senators, that it were granted us to pay rewards to all our citizens! Yet what we have promised we will repay zealously, and heaped up. But that, I hope, remains for the victors, to whom the Senate's faith will be made good; and since they followed it in the commonwealth's most difficult hour, they shall never have cause to repent their choice. But it is easy to deal well with men by whom we feel ourselves pressed even when they are silent: this is the more admirable, the greater, and most of all the mark of a wise Senate --- to attend with grateful memory the valour of those who have poured out their lives for the fatherland.

\ciceroSection{31} For their honour, would that more came into my mind! Two at least I will not pass over, the two that most present themselves: the one touches the everlasting glory of those bravest men, the other the easing of their kindred's grief and mourning. I am resolved, then, senators, that for the soldiers of the Martian legion, and for those who fell fighting at their side, a monument be raised as ample as may be. Great and past belief are this legion's services to the commonwealth. This was the first to tear itself from Antony's banditry; this held Alba; this betook itself to Caesar; and the Fourth legion, in imitation of it, won an equal glory of valour. The Fourth, victorious, misses no man; of the Martian, a few fell in the very moment of victory. O fortunate death, which, owed to nature, was paid above all for the fatherland!

\ciceroSection{32} But you I judge born for the fatherland --- you whose very name is from Mars, so that the same god seems to have begotten this city for the nations, and you for this city. In flight death is foul; in victory, glorious. For Mars himself is wont to claim as his own pledge the bravest from the battle-line. Those impious men, then, whom you killed will pay, even among the shades below, the penalty of their parricide; but you, who poured out your last breath in victory, have won the seat and place of the blessed. Brief is the life nature has given us; but the memory of a life well rendered up is everlasting. And if that memory were no longer than this life, who would be so deranged as to strive through the greatest toils and perils for the highest praise and glory?

\ciceroSection{33} Gloriously, then, has it gone with you, bravest of soldiers while you lived, but now most hallowed too, in that your valour can be buried neither by the forgetfulness of those now living nor by the silence of those to come, since the Senate and people of Rome will have raised you an immortal monument almost with their own hands. Many armies, often, in the Punic, the Gallic, the Italian wars, were famous and great, yet to none was such a kind of honour given. And would that we could give greater, since from you we have received the greatest! You turned the raging Antony away from the city; you drove him back as he strove to return. There shall stand, then, a structure reared in magnificent work, and letters carved upon it, everlasting witnesses of a godlike valour; and never shall the most grateful speech about you fall silent on the lips of those who either see your monument or hear of it. So, in exchange for the mortal condition of life, you have won immortality.

\ciceroSection{34} But since, senators, the gift of glory is paid to the best and bravest citizens by the honour of a monument, let us console their nearest, for whom this indeed is the best consolation: to the parents, that they have begotten such bulwarks of the commonwealth; to the children, that they will have at home examples of valour; to the wives, that they will be bereft of husbands whom it will be better to praise than to mourn; to the brothers, that they will trust there is in themselves a likeness of valour, as there is of body. And would that by our motions and decrees we could wipe away the tears of all these, or that some such speech could be addressed to them in public, whereby they might lay aside their grief and mourning and rejoice rather that, while many and various kinds of death hang over men, the fairest kind has fallen to their own --- and that these lie neither unburied nor deserted (a thing which, even so, is not held pitiable when it is for the fatherland), nor burned with lowly burial in scattered graves, but covered over by public works and public gifts, by that structure which shall be, unto the memory of eternity, an altar of valour.

\ciceroSection{35} Wherefore it will be the greatest solace of their kin that by one and the same monument there is declared the valour of their own, the devotion of the Roman people, the faith of the Senate, and the memory of a most cruel war --- a war in which, had not so great a valour of the soldiers stood forth, the name of the Roman people would have perished by the parricide of Marcus Antonius. And I move also, senators, that the rewards we promised to give the soldiers when the commonwealth was recovered be paid in full and heaped up, when the time comes, to the living and the victorious; but as for those among them, to whom those rewards were promised, who have died for the fatherland, I move that the same be given to their parents, their children, their wives, and their brothers.

\ciceroSection{36} But, to embrace my motion at last, I move thus. Whereas Gaius Pansa, consul and imperator, made the beginning of the engagement with the enemy, in which battle the Martian legion, with admirable and incredible valour, defended the liberty of the Roman people, and the legions of recruits did the same; and whereas Gaius Pansa himself, consul and imperator, while he bore himself amid the very midst of the enemy's weapons, received wounds; and whereas Aulus Hirtius, consul and imperator, on hearing of the battle and learning how matters stood, with the bravest and most excellent spirit led his army out of the camp and made an assault upon Marcus Antonius and the enemy's army and slaughtered his forces with utter slaughter, his own army so unharmed that he did not lose even a single soldier; and whereas Gaius ---

\ciceroSection{37} Caesar, propraetor and imperator, by his counsel and diligence successfully defended the camp and routed and slew the enemy's forces that had advanced upon it: for these things the Senate deems and judges that by the valour, the command, the counsel, the gravity, the constancy, the greatness of spirit, and the good fortune of those three imperatores the Roman people has been delivered from a most foul and cruel slavery; and since they have preserved the commonwealth, the city, the temples of the immortal gods, and the goods and fortunes and children of all men, by their struggle and at the peril of their own lives --- that for these things well, bravely, and successfully done, Gaius Pansa and Aulus Hirtius, consuls and imperatores, the one or both, or, if they are absent, Marcus Cornutus the urban praetor, shall appoint a thanksgiving of fifty days at all the sacred couches of the gods.

\ciceroSection{38} And since the valour of the legions has shown itself worthy of those most illustrious imperatores, the Senate will, with the utmost zeal, when the commonwealth is recovered, pay in full what it has heretofore promised to our legions and armies; and since the Martian legion was first to close with the enemy, and so contended against a greater number that they cut down very many while a few of their own fell, and since without any drawing back they poured out their lives for the fatherland; and since with like valour the soldiers of the remaining legions met death for the safety and liberty of the Roman people --- the Senate resolves that Gaius Pansa and Aulus Hirtius, consuls and imperatores, the one or both, if it seem good to them, shall take care that for those who shed their blood for the life, the liberty, and the fortunes of the Roman people, for the city and the temples of the immortal gods, a monument be contracted for and built as ample as may be; and that they bid the urban quaestors to furnish, assign, and pay money to that end, so that it may bear witness, unto the everlasting memory of posterity, to the wickedness of the most cruel enemies and the godlike valour of the soldiers; and that the rewards the Senate has heretofore established for the soldiers be paid to the parents, children, wives, and brothers of those who in this war have died for the fatherland; and that there be given to them what ought to have been given to the soldiers themselves, had those conquered alive who conquered by death.
